\documentclass[12pt,a4paper]{article}

% -------------------------------------------------
% PACOTES BÁSICOS
% -------------------------------------------------
\usepackage[utf8]{inputenc}
\usepackage[T1]{fontenc}
\usepackage[brazil]{babel}
\usepackage{lmodern}
\usepackage{microtype}
\usepackage{geometry}
\usepackage{setspace}
\usepackage{parskip}
\usepackage{xcolor}
\usepackage{graphicx}
\usepackage{booktabs}
\usepackage{tabularx}
\usepackage{longtable}
\usepackage{array}
\usepackage{enumitem}
\usepackage{hyperref}
\usepackage{fancyhdr}
\usepackage{titlesec}
\usepackage{tikz}
\usepackage{float}
\usepackage{seqsplit}
\usetikzlibrary{arrows.meta, positioning, shapes.geometric, shapes.misc}

% -------------------------------------------------
% CONFIGURAÇÃO DA PÁGINA
% -------------------------------------------------
\geometry{left=2.5cm,right=2.5cm,top=2.5cm,bottom=2.5cm}
\onehalfspacing
\hypersetup{
    colorlinks=true,
    linkcolor=blue!45!black,
    urlcolor=blue!45!black,
    citecolor=blue!45!black
}

% -------------------------------------------------
% CORES E ESTILO
% -------------------------------------------------
\definecolor{azulinst}{HTML}{1F4E79}
\definecolor{cinzaclaro}{HTML}{F2F4F7}
\definecolor{cinzamedio}{HTML}{D9E2EC}
\definecolor{verdeproc}{HTML}{2E7D32}
\definecolor{laranja}{HTML}{EF6C00}
\definecolor{vermelho}{HTML}{B00020}

\titleformat{\section}
  {\Large\bfseries\color{azulinst}}
  {\thesection}{0.8em}{}

\titleformat{\subsection}
  {\large\bfseries\color{azulinst}}
  {\thesubsection}{0.8em}{}

\pagestyle{fancy}
\fancyhf{}
\lhead{Relatório do Fluxo de Criação de Cursos}
\rhead{\thepage}
\renewcommand{\headrulewidth}{0.4pt}

\newcolumntype{Y}{>{\raggedright\arraybackslash}X}

% -------------------------------------------------
% COMANDOS ÚTEIS
% -------------------------------------------------
\newcommand{\campo}[1]{\textcolor{azulinst}{\textbf{#1}}}

\newenvironment{caixa}[1]
{%
\vspace{0.4cm}
\noindent
\begin{tikzpicture}
\node[draw=azulinst, fill=cinzaclaro, rounded corners, inner sep=10pt, text width=0.95\textwidth] (box)
\bgroup
\textbf{#1}\par\medskip
}
{%
\egroup;
\end{tikzpicture}
\vspace{0.2cm}
}

% -------------------------------------------------
% DADOS DO DOCUMENTO
% -------------------------------------------------
\title{\textbf{Relatório do Fluxo de Criação de Cursos}\\
\large Processo de conversão, geração, revisão e reserialização para Articulate Rise 360}
\author{\textbf{Responsável:} João e Marina}
\date{\textbf{Versão:} 1.0 \\ \textbf{Data:} \today}

\begin{document}

\maketitle
\thispagestyle{empty}

\begin{center}
\vfill
\begin{tabular}{ll}
\toprule
\textbf{Projeto} & Criação de cursos no Articulate Rise 360 \\
\textbf{Área demandante} & Inserir área \\
\textbf{Responsável técnico} & Inserir responsável \\
\textbf{Ferramentas principais} & Markdown, Articulate Rise 360, SCORM, JSON, Claude Code \\
\textbf{Status} & Em elaboração / Em validação / Concluído \\
\bottomrule
\end{tabular}
\vfill
\end{center}

\newpage
\tableofcontents
\newpage

% -------------------------------------------------
\section{Objetivo do Relatório}
% -------------------------------------------------

Este relatório tem por objetivo documentar a \textit{pipeline} de criação de cursos digitais assistida por modelos de Inteligência Artificial (Large Language Models -- LLMs), descrevendo as etapas, os artefatos produzidos e os mecanismos de controle adotados ao longo do processo.

A \textit{pipeline} inicia-se com a conversão de um plano de aula em um roteiro de conteúdo estruturado por um LLM. Em seguida, o roteiro é utilizado pelo Articulate Rise 360 AI, juntamente com um prompt de instruções, para construir a primeira versão do curso utilizando os componentes da plataforma. Posteriormente, o pacote SCORM é exportado, desserializado e submetido a uma revisão assistida pelo Claude (Anthropic), responsável pelo refinamento do conteúdo e pela preservação da estrutura do projeto.

Este documento também apresenta os pontos de verificação, os ciclos de retrabalho e os procedimentos adotados para garantir a consistência do curso durante todo o fluxo de produção.

A documentação desta \textit{pipeline} destina-se a servir como referência para padronização, reprodução e evolução do processo de criação de cursos da ArcTel.


% -------------------------------------------------
\section{Visão Geral do Fluxo}

\begin{itemize}[leftmargin=*]
    \item Conversão do plano de aula em roteiro instrucional no formato Markdown.
    \item Geração da versão inicial do curso no Articulate Rise 360.
    \item Verificação visual inicial do curso.
    \item Exportação do pacote SCORM.
    \item Desserialização do pacote SCORM em arquivos JSON.
    \item Separação do conteúdo em lições para revisão técnica.
    \item Análise assistida da estrutura de blocos por LLM.
    \item Revisão e edição do conteúdo e da estrutura de blocos.
    \item Reserialização do projeto e geração do pacote SCORM final.
\end{itemize}
% -------------------------------------------------

\begin{figure}[h!]
\centering
\resizebox{0.95\textwidth}{!}{%
\begin{tikzpicture}[
    node distance=1.05cm,
    etapa/.style={rectangle, rounded corners, draw=azulinst, fill=cinzaclaro, align=center, text width=4.2cm, minimum height=0.9cm},
    decisao/.style={diamond, draw=laranja, fill=orange!10, align=center, aspect=2.2, text width=3.2cm, inner sep=1pt},
    fim/.style={rectangle, rounded corners, draw=verdeproc, fill=green!10, align=center, text width=3.5cm, minimum height=0.9cm},
    seta/.style={-{Latex[length=3mm]}, thick, color=azulinst}
]

\node[etapa] (e1) {1. Converter plano de aula para roteiro};
\node[etapa, below=of e1] (e2) {2. Carregar prompt e roteiro no Articulate};
\node[etapa, below=of e2] (e3) {3. Gerar curso};
\node[etapa, below=of e3] (e4) {4. Verificação visual mínima};
\node[etapa, below=of e4] (e5) {5. Gerar SCORM};
\node[etapa, below=of e5] (e6) {6. Desserializar e separar lições em JSON};
\node[decisao, below=of e6] (d1) {7. Estrutura de blocos adequada?};
\node[etapa, right=2.2cm of d1] (e9) {9. Alterar conteúdo sem quebrar estrutura};
\node[etapa, below=of d1] (e10) {10. Reserializar e gerar novo SCORM};
\node[fim, below=of e10] (fim) {11. Fim};

\draw[seta] (e1) -- (e2);
\draw[seta] (e2) -- (e3);
\draw[seta] (e3) -- (e4);
\draw[seta] (e4) -- (e5);
\draw[seta] (e5) -- (e6);
\draw[seta] (e6) -- (d1);
\draw[seta] (d1) -- node[right]{Sim} (e10);
\draw[seta] (e10) -- (fim);
\draw[seta] (d1) -- node[above]{Não / ajustes} (e9);
\draw[seta] (e9) |- (e5);

\end{tikzpicture}}
\caption{Fluxo macro de criação, revisão e reserialização de cursos.}
\end{figure}

\subsection{Etapa 1 -- Conversão do plano de aula para roteiro}

A primeira etapa da \textit{pipeline} consiste na transformação do plano de aula em um roteiro instrucional estruturado para processamento pelo Articulate Rise 360.

O plano de aula constitui o documento-base do curso, contendo a estrutura pedagógica a ser desenvolvida, incluindo carga horária, público-alvo, objetivos de aprendizagem, organização das oficinas, conteúdos, atividades práticas, formas de avaliação e demais orientações instrucionais.

A conversão é realizada por um modelo de linguagem (LLM), utilizando um \textit{prompt} desenvolvido especificamente para essa finalidade. Esse \textit{prompt} foi construído a partir de técnicas de engenharia de prompts (\textit{prompt engineering}) e refinado iterativamente por meio de sucessivos ciclos de experimentação. Em cada iteração, foram avaliadas diferentes estratégias de instrução com o objetivo de identificar quais produziam melhores resultados quando processadas pela Inteligência Artificial do Articulate Rise 360. O resultado desse processo foi um conjunto de instruções capaz de aumentar a aderência entre o curso gerado e a estrutura pedagógica originalmente definida, explorando de forma mais eficiente os componentes e recursos disponibilizados pela plataforma. 

Como produto dessa etapa, é gerado um roteiro em formato Markdown que reúne dois elementos distintos: o conteúdo instrucional do curso e um conjunto de instruções de montagem destinadas exclusivamente à IA do Articulate Rise 360. Essas instruções, denominadas \textit{cues} de montagem, orientam a ferramenta quanto à organização das lições, seleção dos blocos mais adequados, preservação de elementos específicos e utilização dos recursos disponíveis na plataforma. Embora esse mecanismo aumente significativamente a qualidade dos cursos produzidos, a IA do Articulate Rise 360 não segue integralmente todas as instruções fornecidas, tornando necessárias etapas posteriores de revisão e refinamento da estrutura e do conteúdo.

\begin{longtable}{p{0.25\textwidth}p{0.68\textwidth}}
\toprule
\campo{Entrada} &
Arquivo \texttt{plano\_de\_aula\_xxx.md}, preferencialmente elaborado no formato \textit{Essentials/Hands-on}. \\
\vspace{0.05cm}
\campo{Ferramentas} &
\vspace{0.05cm}
LLM de propósito geral e \textit{prompt} especializado para conversão do plano de aula em roteiro instrucional estruturado. \\
\vspace{0.05cm}
\campo{Ação} &
\vspace{0.05cm}
Aplicar o \textit{prompt} de conversão para transformar o plano de aula em um roteiro instrucional compatível com a IA do Articulate Rise 360. \\
\vspace{0.05cm}
\campo{Objetivo} &
\vspace{0.05cm}
Produzir um roteiro contendo o conteúdo completo do curso e as instruções auxiliares necessárias para orientar a montagem automática no Articulate Rise 360. \\
\vspace{0.05cm}
\campo{Saída} &
\vspace{0.05cm}
Arquivo \texttt{ROTEIRO\_xxx.md}. \\
\vspace{0.05cm}
\campo{Controle mínimo} &
\vspace{0.05cm}
Verificar se os objetivos de aprendizagem, oficinas, atividades práticas, entregáveis e \textit{cues} de montagem foram corretamente gerados e preservados. \\

\bottomrule
\end{longtable}

\subsection{Etapa 2 -- Geração inicial do curso no Articulate Rise 360}

Concluída a geração do roteiro instrucional, inicia-se a etapa de criação do curso no Articulate Rise 360. Nessa fase, o roteiro produzido na etapa anterior é fornecido à Inteligência Artificial da plataforma juntamente com um segundo \textit{prompt}, desenvolvido especificamente para orientar a IA quanto à forma de interpretar o roteiro e utilizar os componentes disponíveis no Rise 360.

Diferentemente do \textit{prompt} utilizado na etapa anterior, cujo objetivo é converter o plano de aula em um roteiro estruturado, o \textit{prompt} do Articulate tem como finalidade orientar a montagem do curso propriamente dita. Suas instruções procuram induzir a IA a selecionar os blocos mais adequados para cada tipo de conteúdo, preservar a organização das lições, respeitar as \textit{cues} de montagem presentes no roteiro e manter, na maior medida possível, a estrutura didática originalmente planejada.

Ao término dessa etapa, obtém-se a primeira versão funcional do curso no ambiente do Articulate Rise 360. Essa versão ainda não é considerada definitiva, pois a IA da plataforma pode interpretar parcialmente algumas instruções ou reorganizar determinados elementos do roteiro, justificando as etapas posteriores de revisão e refinamento da \textit{pipeline}.

\begin{longtable}{p{0.25\textwidth}p{0.68\textwidth}}
\toprule
\campo{Entrada} &
Arquivo \texttt{ROTEIRO\_xxx.md} e \textit{prompt} de geração do curso no Articulate Rise 360. \\

\campo{Ferramenta} &
Articulate Rise 360 AI. \\

\campo{Ação} &
Inserir o roteiro e o \textit{prompt} no ambiente de geração do Articulate Rise 360 e executar a criação automática do curso. \\

\campo{Objetivo} &
Produzir a primeira versão do curso utilizando os blocos, componentes e recursos disponibilizados pela plataforma. \\

\campo{Saída} &
Primeira versão do curso disponível para inspeção no ambiente do Articulate Rise 360. \\

\campo{Controle mínimo} &
Verificar se todas as lições foram criadas, se a organização geral do curso foi preservada e se não ocorreram omissões ou reorganizações significativas do conteúdo. \\

\bottomrule
\end{longtable}

\subsubsection*{Procedimento de criação no Articulate Rise 360}

As Figuras a seguir ilustram o procedimento de criação do curso no Articulate Rise 360.

\begin{figure}[H]
    \centering
    \includegraphics[width=\textwidth]{figuras/Figura2.png}
    \caption{Tela inicial do Articulate Rise 360 e seleção da opção para criação de um novo curso utilizando Inteligência Artificial (Create with AI).}
    \label{fig:rise01}
\end{figure}

\begin{figure}[H]
    \centering
    \includegraphics[scale=0.5]{figuras/Figura3.png}
    %\fbox{\rule{0pt}{7cm}\rule{0.95\textwidth}{0pt}}
    \caption{Inserção do \textit{prompt} de geração do curso e do roteiro instrucional produzido na Etapa 1. Para prosseguir com o curso clique "Genarate course details".}
    \label{fig:rise02}
\end{figure}

\begin{figure}[H]
    \centering
    \includegraphics[scale=0.4]{figuras/Figura4.1.png}
    \includegraphics[scale=0.4]{figuras/Figura4.2.png}
    \includegraphics[scale=0.4]{figuras/Figura4.3.png}
    %\fbox{\rule{0pt}{7cm}\rule{0.95\textwidth}{0pt}}
    \caption{Configuração dos parâmetros do curso, com destaque para os objetivos de aprendizagem, carga horária e demais informações apresentadas pelo Articulate Rise 360 antes da geração do conteúdo.
    É importante ressaltar que esses parâmetros são gerados automaticamente e para prosseguir com a criação do curso clique "Generate course outline".}
    \label{fig:rise03}
\end{figure}

\begin{figure}[H]
    \centering
    \includegraphics[scale=0.5]{figuras/Figura5.png}
    \caption{Configuração das lições do curso geradas automaticamente pela AI do Articulate Rise 360. É importante verificar se a quantidade e os temas das lições estão de acordo com o roteiro produzido na Etapa 1. Clique em "Course content options".}
    \label{fig:rise04}
\end{figure}

\begin{figure}[H]
    \centering
    %\includegraphics[width=\textwidth]{figuras/Figura3.png}
    \includegraphics[scale=0.5]{figuras/Figura6.png}
    %\fbox{\rule{0pt}{7cm}\rule{0.95\textwidth}{0pt}}
    \caption{Clique em "Create course draft" para finalizar a criação da primeira versão do curso.}
    \label{fig:rise05}
\end{figure}

\begin{figure}[H]
    \centering
    \includegraphics[scale=0.5]{figuras/Figura8.1.png}
    \includegraphics[scale=0.5]{figuras/Figura8.2.png}
    \caption{Primeira versão do curso gerada automaticamente no ambiente do Articulate Rise 360, para editar manualmente as lições clique em "Edit Content".}
    \label{fig:rise06}
\end{figure}

\begin{figure}[H]
    \centering
    \includegraphics[scale=0.4]{figuras/Figura9.1.png}
    \includegraphics[scale=0.4]{figuras/Figura9.2.png}
    \caption{Primeira versão do curso gerada automaticamente no ambiente do Articulate Rise 360, para editar manualmente as lições clique em "Edit Content".}
    \label{fig:rise06}
\end{figure}

\begin{figure}[H]
    \centering
    %\includegraphics[width=\textwidth]{figuras/Figura3.png}
    \includegraphics[scale=0.4]{figuras/Figura7.png}
    %\fbox{\rule{0pt}{7cm}\rule{0.95\textwidth}{0pt}}
    \caption{Primeira versão do curso gerada automaticamente no ambiente do Articulate Rise 360, para editar manualmente as lições clique em "Edit Content".}
    \label{fig:rise06}
\end{figure}

\subsection{Etapa 3 -- Verificação visual inicial}

Após a geração do curso no Articulate Rise 360, realiza-se uma inspeção visual inicial com o objetivo de verificar se a Inteligência Artificial da plataforma interpretou corretamente o roteiro instrucional e as instruções de montagem fornecidas.

Essa etapa não tem como finalidade revisar o conteúdo pedagógico em profundidade, mas identificar erros evidentes decorrentes do processo de geração automática. Durante o desenvolvimento desta \textit{pipeline}, observou-se que a IA do Articulate Rise 360 pode, ocasionalmente, produzir resultados inconsistentes, como interpretar o \textit{prompt} de instruções como parte do conteúdo do curso, gerar o curso em idioma diferente do especificado, omitir seções do roteiro, reorganizar lições de maneira inadequada ou utilizar blocos incompatíveis com as instruções fornecidas.

A identificação precoce dessas inconsistências evita que erros sejam propagados para as etapas subsequentes da \textit{pipeline}, reduzindo retrabalho e garantindo que apenas versões minimamente consistentes sejam exportadas para o pacote SCORM.

\begin{longtable}{p{0.25\textwidth}p{0.68\textwidth}}
\toprule

\campo{Entrada} &
Primeira versão do curso gerada no Articulate Rise 360. \\[0.5cm]

\campo{Ferramenta} &
Articulate Rise 360. \\[0.5cm]

\campo{Ação} &
Realizar inspeção visual da estrutura geral do curso e verificar a consistência da geração automática. \\[1cm]

\campo{Objetivo} &
Identificar inconsistências produzidas pela IA antes da exportação do pacote SCORM. \\[1cm]

\campo{Saída} &
Curso aprovado para exportação ou indicação de necessidade de nova geração do curso. \\[1cm]

\campo{Controle mínimo} &
Verificar, no mínimo:
\begin{itemize}[nosep,leftmargin=*]
    \item idioma do curso;
    \item título e descrição do curso;
    \item quantidade e ordem das lições;
    \item correspondência entre o roteiro e o conteúdo gerado;
    \item utilização adequada dos blocos do Rise 360;
    \item existência de atividades, imagens e elementos interativos previstos;
    \item ausência de trechos do \textit{prompt} ou de instruções internas exibidos ao aluno.
\end{itemize}
\\[1cm]

\bottomrule
\end{longtable}

\subsection{Etapa 4 -- Geração do SCORM}

Após a validação inicial do curso no Articulate Rise 360, realiza-se a exportação do projeto em formato SCORM. Essa etapa transforma o curso criado na plataforma em um pacote técnico no formato \texttt{.zip}, compatível com ambientes virtuais de aprendizagem (LMS) e também manipulável nas etapas posteriores da \textit{pipeline}.

A exportação deve seguir um procedimento padronizado, especialmente quanto à escolha do padrão SCORM, à identificação da versão do curso e ao armazenamento do arquivo exportado em uma pasta específica do projeto. Esse cuidado permite rastrear as versões geradas, evitar substituições indevidas e garantir que o pacote correto seja utilizado nas etapas de desserialização e revisão.

\begin{longtable}{p{0.25\textwidth}p{0.68\textwidth}}
\toprule

\campo{Entrada} &
Curso aprovado na validação inicial. \\[0.5cm]

\campo{Ferramenta} &
Articulate Rise 360. \\[0.5cm]

\campo{Ação} &
Exportar o curso como pacote SCORM, selecionando o padrão adequado e salvando o arquivo em pasta específica do projeto. \\[1cm]

\campo{Objetivo} &
Produzir um pacote técnico publicável em LMS e utilizável nas etapas posteriores de desserialização, análise e revisão. \\[1cm]

\campo{Saída} &
Arquivo \texttt{.zip} contendo o pacote SCORM. \\[0.5cm]

\campo{Controle mínimo} &
Registrar nome do curso, versão do SCORM, data da exportação, local de salvamento e versão do pacote exportado. \\[0.05cm]

\bottomrule
\end{longtable}

\subsubsection*{Procedimento de exportação do pacote SCORM}

As figuras a seguir devem documentar os passos de exportação do curso no Articulate Rise 360.

\begin{figure}[H]
    \centering
    \includegraphics[scale=0.4]{figuras/figura10.png}
    \caption{Acesso à opção de exportação do curso no Articulate Rise 360 e selecione a opção de exportação para LMS. (Publish --> LMS)}
    \label{fig:scorm01}
\end{figure}

\begin{figure}[H]
    \centering
     \includegraphics[scale=0.4]{figuras/figura12.png}
    \caption{Escolha da versão do padrão SCORM a ser utilizada na exportação do curso (SCORM 1.2), clique em Download (canto superior direito).}
    \label{fig:scorm03}
\end{figure}

\begin{figure}[H]
    \centering
     \includegraphics[scale=0.4]{figuras/figura13.png}
    \caption{Por padrão, o salvamento do pacote SCORM exportado é feito por meio de uma pasta zip nos downloads, salve em uma pasta específica do projeto.}
    \label{fig:scorm04}
\end{figure}

\subsection{Etapa 6 -- Desserialização e separação em lições}

Após a exportação do curso em formato SCORM, o pacote \texttt{.zip} deve ser extraído para permitir o acesso aos arquivos internos gerados pelo Articulate Rise 360. O conteúdo principal do curso não é armazenado em arquivos JSON independentes, mas encontra-se incorporado ao arquivo \texttt{index.html}, localizado na pasta \texttt{scormcontent}.

No interior desse arquivo existe uma função JavaScript responsável por carregar o curso. Essa função possui formato semelhante ao apresentado a seguir:

\begin{verbatim}
async function __fetchCourse() {
    return Promise.resolve(deserialize("CONTEUDO_SERIALIZADO"))
}
\end{verbatim}

O argumento da função \texttt{deserialize()} consiste em uma longa sequência de caracteres codificada (Base64 ou formato equivalente utilizado pelo Articulate Rise 360), que representa todo o curso serializado. Essa sequência contém a estrutura completa do curso, incluindo lições, blocos, textos, imagens, interações, metadados e demais informações necessárias para reconstrução do projeto.

Atualmente, a localização dessa sequência é realizada manualmente pelo operador. Após localizar a função \texttt{\_\_fetchCourse()}, todo o conteúdo serializado deve ser copiado e inserido em um dos scripts desenvolvidos em Python para esta \textit{pipeline}. Esses scripts têm como única finalidade desserializar o conteúdo do curso e convertê-lo para arquivos JSON legíveis, preservando integralmente sua estrutura original.

Como resultado desse processo, o curso passa a ser representado por um conjunto de arquivos JSON, separados por lição. Essa representação torna o conteúdo acessível para inspeção técnica e edição controlada, sem necessidade de manipular diretamente o conteúdo serializado presente no arquivo \texttt{index.html}.

Os arquivos JSON produzidos nessa etapa constituem a entrada para a etapa seguinte da \textit{pipeline}. Nela, uma LLM — atualmente o Claude Code — recebe as lições juntamente com um \textit{prompt} de revisão desenvolvido especificamente para esse processo. O objetivo dessa revisão é comparar a estrutura e o conteúdo do curso gerado com o roteiro instrucional original, identificar inconsistências produzidas pela IA do Articulate Rise 360 e propor melhorias no conteúdo textual, preservando rigorosamente toda a estrutura técnica do JSON. O \textit{prompt} empregado estabelece restrições explícitas para impedir alterações em identificadores, tipos de blocos, metadados, posições, arrays e demais elementos estruturais do curso. :contentReference[oaicite:0]{index=0}

\begin{longtable}{p{0.25\textwidth}p{0.68\textwidth}}
\toprule

\campo{Entrada} &
Pacote SCORM exportado no formato \texttt{.zip}. \\[0.7cm]

\campo{Ferramentas} &
Editor de texto, scripts em Python para desserialização e Claude Code para revisão assistida. \\[1.4cm]

\campo{Ação} &
Extrair o pacote SCORM, localizar manualmente a função \texttt{\_\_fetchCourse()}, copiar o conteúdo serializado, desserializá-lo utilizando os scripts em Python e separar o curso em arquivos JSON organizados por lição. \\[2.5cm]

\campo{Objetivo} &
Transformar o conteúdo serializado do curso em arquivos JSON legíveis e preparados para análise e revisão assistida por LLM. \\[2cm]

\campo{Saída} &
Arquivos JSON separados por lição e prontos para processamento pela LLM. \\[1.4cm]

\campo{Controle mínimo} &
Verificar se todas as lições foram corretamente desserializadas, se os arquivos JSON permanecem válidos e se a estrutura do curso foi integralmente preservada. \\[0.7cm]

\bottomrule
\end{longtable}

\subsubsection*{Procedimento de desserialização do curso}

As figuras a seguir ilustram as etapas necessárias para obtenção dos arquivos JSON utilizados nas etapas posteriores da \textit{pipeline}.

\begin{figure}[H]
    \centering
    \includegraphics[scale=0.4]{figuras/figura14.png}
    \caption{Estrutura de diretórios do pacote SCORM após a extração do arquivo \texttt{.zip}, destacando a pasta \texttt{scormcontent} com o arquivo \texttt{index}.}
    \label{fig:desserializacao01}
\end{figure}

\begin{figure}[H]
    \centering
    \includegraphics[scale=0.4]{figuras/figura15.png}
    \caption{Identificação da função \texttt{\_\_fetchCourse()} e do conteúdo serializado utilizado pelo Articulate Rise 360 para reconstrução do curso. Use \texttt{Ctrl + F} para buscar a função no código. }
    \label{fig:desserializacao03}
\end{figure}

\begin{figure}[H]
    \centering
    \includegraphics[scale=0.4]{figuras/Figura16.png}
    \caption{Identificação do arquivo \texttt{desserialize.py} na estrutura de páginas do script python.}
    \label{fig:desserializacao04}
\end{figure}

\begin{figure}[H]
    \centering
    \includegraphics[scale=0.4]{figuras/figura17.png}
    \caption{Inserção do conteúdo serializado no script Python responsável pela desserialização do curso.}
    \label{fig:desserializacao04}
\end{figure}

\begin{figure}[H]
    \centering
    \includegraphics[scale=0.4]{figuras/figura18.png}
    \includegraphics[scale=0.4]{figuras/figura19.png}
    \caption{Atenção a sintaxe, é importante retirar os parênteses no final da linha e mudar os nomes para se adequar ao curso a ser desserializado. Depois já está pronto para rodar.}
    \label{fig:desserializacao05}
\end{figure}

\begin{figure}[H]
    \centering
    \includegraphics[scale=0.4]{figuras/figura20.png}
    \includegraphics[scale=0.4]{figuras/figura21.png}
    \caption{Arquivos JSON produzidos após a desserialização, é importante movê-los para dentro da pasta \texttt{cursos_-desserializados}.}
    \label{fig:desserializacao05}
\end{figure}

\begin{figure}[H]
    \centering
    \includegraphics[scale=0.4]{figuras/figura18.png}
    \caption{Arquivos JSON produzidos após a desserialização, organizados individualmente por lição e preparados para revisão assistida por LLM.}
    \label{fig:desserializacao05}
\end{figure}
\subsection{Etapa 7 -- Avaliação da estrutura de blocos do Articulate}

\begin{longtable}{p{0.25\textwidth}p{0.68\textwidth}}
\toprule
\campo{Entrada} & Arquivos JSON extraídos do SCORM. \\
\campo{Ação} & Avaliar alterações na estrutura de blocos do Articulate. \\
\campo{Objetivo} & Verificar se os blocos permanecem tecnicamente íntegros e compatíveis com o Rise 360. \\
\campo{Saída} & Diagnóstico da estrutura: aprovada, aprovada com ressalvas ou reprovada. \\
\campo{Controle mínimo} & Não alterar identificadores, tipos de bloco, dependências, metadados ou campos estruturais sensíveis. \\
\bottomrule
\end{longtable}

\subsection{Etapa 8 -- Retorno para nova geração de SCORM}

Caso a estrutura seja considerada adequada, o fluxo retorna à etapa de geração do SCORM, permitindo nova exportação a partir da versão ajustada. Caso sejam detectados problemas relevantes, deve-se registrar a inconsistência e retornar à etapa de ajuste controlado.

\subsection{Etapa 9 -- Revisão assistida do conteúdo e edição segura do JSON}

Nesta etapa é realizada a revisão assistida do curso utilizando uma LLM. Até o momento, o processo foi desenvolvido e validado utilizando o Claude Code, embora possa ser adaptado para outros modelos de linguagem que possuam capacidade de compreender documentos extensos e editar arquivos mantendo sua estrutura.

A revisão é realizada utilizando dois artefatos produzidos nas etapas anteriores da \textit{pipeline}: o roteiro instrucional gerado na Etapa 1 e os arquivos JSON produzidos na Etapa 6. O roteiro funciona como referência pedagógica, enquanto os arquivos JSON representam a implementação efetivamente produzida pelo Articulate Rise 360.

O objetivo da revisão é comparar o curso gerado com o roteiro original, identificar divergências introduzidas durante a geração automática e melhorar o conteúdo textual das lições sem alterar a arquitetura técnica do curso. Para isso é utilizado um \textit{prompt} específico de revisão, desenvolvido para orientar a LLM a editar exclusivamente os campos textuais do JSON, preservando rigorosamente todos os elementos estruturais do projeto.

Esse \textit{prompt} estabelece diversas restrições de segurança estrutural, impedindo alterações em identificadores, tipos de bloco, metadados, posições, arrays, configurações internas e demais informações utilizadas pelo Articulate Rise 360 para reconstrução do curso. Dessa forma, a revisão concentra-se apenas na melhoria da qualidade pedagógica do conteúdo, sem comprometer o funcionamento da estrutura técnica posteriormente utilizada na reserialização do curso.

\begin{longtable}{p{0.25\textwidth}p{0.68\textwidth}}
\toprule

\campo{Entrada} &
Roteiro instrucional produzido na Etapa 1, arquivos JSON separados por lição e \textit{prompt} de revisão assistida. \\[0.05cm]

\campo{Ferramenta} &
Claude Code (LLM) e \textit{prompt} de revisão estruturada. \\[0.05cm]

\campo{Ação} &
Comparar o conteúdo dos arquivos JSON com o roteiro original e revisar exclusivamente os campos textuais permitidos, preservando integralmente a estrutura técnica do curso. \\[0.05cm]

\campo{Objetivo} &
Melhorar a clareza, profundidade, qualidade instrucional e fidelidade pedagógica do curso, sem modificar a arquitetura técnica do JSON. \\[0.05cm]

\campo{Saída} &
Arquivos JSON revisados, mantendo compatibilidade integral com o Articulate Rise 360. \\[0.05cm]

\campo{Controle mínimo} &
Confirmar que não houve alteração em IDs, tipos de bloco, posições, metadados, arrays, configurações internas ou qualquer outro elemento estrutural do curso. As alterações devem restringir-se aos campos textuais autorizados pelo \textit{prompt} de revisão. \\[0.05cm]

\bottomrule
\end{longtable}

\subsubsection*{Procedimento de revisão assistida}

As figuras a seguir ilustram o procedimento de revisão dos arquivos JSON utilizando uma LLM.

\begin{figure}[H]
    \centering
    %\includegraphics[width=\textwidth]{figuras/etapa9_01_prompt.png}
    \fbox{\rule{0pt}{7cm}\rule{0.95\textwidth}{0pt}}
    \caption{Prompt de revisão utilizado para orientar a edição segura dos arquivos JSON.}
    \label{fig:revisao01}
\end{figure}

\begin{figure}[H]
    \centering
    %\includegraphics[width=\textwidth]{figuras/etapa9_02_contexto.png}
    \fbox{\rule{0pt}{7cm}\rule{0.95\textwidth}{0pt}}
    \caption{Envio do roteiro instrucional e dos arquivos JSON como contexto para a revisão assistida.}
    \label{fig:revisao02}
\end{figure}

\begin{figure}[H]
    \centering
    %\includegraphics[width=\textwidth]{figuras/etapa9_03_edicao.png}
    \fbox{\rule{0pt}{7cm}\rule{0.95\textwidth}{0pt}}
    \caption{Processo de edição dos campos textuais preservando a estrutura técnica do JSON.}
    \label{fig:revisao03}
\end{figure}

\begin{figure}[H]
    \centering
    %\includegraphics[width=\textwidth]{figuras/etapa9_04_relatorio.png}
    \fbox{\rule{0pt}{7cm}\rule{0.95\textwidth}{0pt}}
    \caption{Relatório gerado pela LLM contendo o resumo das alterações e a validação da preservação da estrutura técnica do curso.}
    \label{fig:revisao04}
\end{figure}

\subsection{Etapa 10 -- Reserialização do curso e validação}

Após a conclusão da revisão assistida, os arquivos JSON modificados devem ser convertidos novamente para o formato serializado utilizado pelo Articulate Rise 360. Essa etapa constitui o processo inverso da desserialização apresentada na Etapa 6.

A reserialização é realizada por meio de scripts desenvolvidos em Python, responsáveis por reconstruir o conteúdo serializado a partir dos arquivos JSON revisados. Como resultado, é produzido um arquivo texto, denominado \texttt{NomeCurso\_base64\_editado.txt}, contendo a nova representação serializada do curso.

Em seguida, o arquivo \texttt{index.html}, localizado na pasta \texttt{scormcontent}, deve ser aberto em um editor de texto. O conteúdo serializado anteriormente presente como argumento da função \texttt{deserialize()} deve ser integralmente substituído pela nova sequência gerada no arquivo \texttt{NomeCurso\_base64\_editado.txt}, mantendo inalterada a estrutura da função \texttt{\_\_fetchCourse()}.

\begin{verbatim}
async function __fetchCourse() {
    return Promise.resolve(deserialize("CONTEUDO_SERIALIZADO_EDITADO"))
}
\end{verbatim}

Ou seja, a substituição deve ocorrer apenas no trecho correspondente a \texttt{CONTEUDO\_SERIALIZADO}, sem alterar o nome da função, a chamada \texttt{Promise.resolve()}, a função \texttt{deserialize()} ou a estrutura JavaScript ao redor.

Após a substituição, o arquivo \texttt{index.html} deve ser salvo e executado em um navegador Web. Caso o processo tenha sido realizado corretamente e nenhuma alteração estrutural indevida tenha sido introduzida durante a revisão dos arquivos JSON, o curso será carregado normalmente, já incorporando todas as modificações realizadas pela LLM.

Essa etapa constitui a principal validação da integridade estrutural do curso. Caso o conteúdo não seja carregado corretamente, isso normalmente indica que algum elemento estrutural do JSON foi alterado inadvertidamente durante a etapa de revisão, exigindo nova inspeção dos arquivos modificados.

\begin{longtable}{p{0.25\textwidth}p{0.68\textwidth}}
\toprule

\campo{Entrada} &
Arquivos JSON revisados e validados. \\[0.05cm]

\campo{Ferramentas} &
Scripts em Python para reserialização, editor de texto e navegador Web. \\[0.05cm]

\campo{Ação} &
Reserializar os arquivos JSON, copiar a nova sequência serializada gerada em \texttt{NomeCurso\_base64\_editado.txt}, substituir o conteúdo presente no argumento da função \texttt{deserialize()} no arquivo \texttt{index.html} e validar o funcionamento do curso em navegador Web. \\[0.05cm]

\campo{Objetivo} &
Reconstruir o curso utilizando o conteúdo revisado e verificar a preservação da estrutura técnica após a edição assistida por LLM. \\[0.05cm]

\campo{Saída} &
Arquivo \texttt{NomeCurso\_base64\_editado.txt} e curso atualizado no arquivo \texttt{index.html}. \\[0.05cm]

\campo{Controle mínimo} &
Verificar se o curso é carregado corretamente no navegador, se todas as lições permanecem acessíveis, se os blocos interativos continuam funcionais e se o conteúdo revisado é exibido corretamente. \\[0.05cm]

\bottomrule
\end{longtable}

\subsubsection*{Procedimento de reserialização}

As figuras a seguir ilustram as etapas necessárias para reconstrução do curso a partir dos arquivos JSON revisados.

\begin{figure}[H]
    \centering
    %\includegraphics[width=\textwidth]{figuras/etapa10_01_python.png}
    \fbox{\rule{0pt}{7cm}\rule{0.95\textwidth}{0pt}}
    \caption{Execução do script Python responsável pela reserialização dos arquivos JSON.}
    \label{fig:reserializacao01}
\end{figure}

\begin{figure}[H]
    \centering
    %\includegraphics[width=\textwidth]{figuras/etapa10_02_base64.png}
    \fbox{\rule{0pt}{7cm}\rule{0.95\textwidth}{0pt}}
    \caption{Arquivo \texttt{NomeCurso\_base64\_editado.txt} produzido após a reserialização.}
    \label{fig:reserializacao02}
\end{figure}

\begin{figure}[H]
    \centering
    %\includegraphics[width=\textwidth]{figuras/etapa10_03_index.png}
    \fbox{\rule{0pt}{7cm}\rule{0.95\textwidth}{0pt}}
    \caption{Substituição do conteúdo serializado no argumento da função \texttt{deserialize()} presente no arquivo \texttt{index.html}.}
    \label{fig:reserializacao03}
\end{figure}

\begin{figure}[H]
    \centering
    %\includegraphics[width=\textwidth]{figuras/etapa10_04_teste.png}
    \fbox{\rule{0pt}{7cm}\rule{0.95\textwidth}{0pt}}
    \caption{Validação do curso atualizado por meio da execução do arquivo \texttt{index.html} em navegador Web.}
    \label{fig:reserializacao04}
\end{figure}

\subsubsection*{Critério de sucesso da reserialização}

Considera-se que a reserialização foi bem-sucedida quando o curso é carregado integralmente no navegador, todas as lições permanecem navegáveis, os blocos interativos continuam funcionais e o conteúdo revisado é exibido corretamente. Esse teste valida simultaneamente a etapa de reserialização e a preservação da estrutura técnica durante a revisão assistida.

\subsection{Etapa 11 -- Encerramento}

O processo é encerrado quando o novo pacote SCORM está funcional, validado e identificado com versão, data e responsável. Recomenda-se manter arquivados o plano de aula original, roteiro, prompts utilizados, SCORM inicial, JSON revisado e SCORM final.

% -------------------------------------------------
\section{Matriz de Artefatos}
% -------------------------------------------------

\begin{table}[H]
\centering
\small
\renewcommand{\arraystretch}{1.25}

\begin{tabularx}{\textwidth}{
>{\ttfamily\raggedright\arraybackslash}p{0.30\textwidth}
>{\raggedright\arraybackslash}X
>{\centering\arraybackslash}p{0.10\textwidth}
}
\toprule
\textnormal{\textbf{Artefato}} &
\textbf{Descrição} &
\textbf{Etapa} \\
\midrule

plano\_de\_aula\_xxx.md &
Plano de aula original contendo a estrutura pedagógica, carga horária, objetivos, conteúdos e organização das oficinas. &
1 \\

ROTEIRO\_xxx.md &
Roteiro instrucional estruturado para processamento pelo Articulate Rise 360, contendo o conteúdo do curso e as \textit{cues} de montagem. &
1--2 \\

promptArticulate.md &
Prompt utilizado para orientar a IA do Articulate Rise 360 durante a geração inicial do curso. &
2 \\

Curso no Rise 360 &
Primeira versão navegável do curso gerada automaticamente pela plataforma. &
2--4 \\

SCORM\_Inicial.zip &
Primeiro pacote SCORM exportado pelo Articulate Rise 360. &
5 \\

Arquivos JSON &
Curso desserializado e organizado em arquivos JSON separados por lição. &
6--9 \\

\texttt{\seqsplit{NomeCurso\_base64\_editado.txt}} &
Conteúdo serializado reconstruído a partir dos arquivos JSON revisados. &
10 \\

SCORM Final &
Versão final do curso validada e pronta para distribuição ou publicação em LMS. &
10 \\

\bottomrule
\end{tabularx}

\caption{Principais artefatos produzidos ao longo da \textit{pipeline} de criação de cursos.}

\end{table}

% -------------------------------------------------
\section{Critérios de Qualidade}
% -------------------------------------------------

\begin{longtable}{p{0.27\textwidth}p{0.63\textwidth}p{0.08\textwidth}}
\toprule
\textbf{Critério} & \textbf{Verificação} & \textbf{OK?} \\
\midrule
Completude do roteiro & Todas as oficinas, atividades e entregáveis previstos no plano foram preservados. & $\square$ \\
Coerência instrucional & O curso apresenta progressão lógica, com teoria apenas quando necessária para a prática. & $\square$ \\
Clareza visual & Títulos, blocos e elementos visuais estão legíveis e bem distribuídos. & $\square$ \\
Integridade técnica & O JSON foi alterado sem rompimento da estrutura interna do Articulate. & $\square$ \\
Funcionalidade do SCORM & O pacote abre, navega e registra progresso no ambiente de teste. & $\square$ \\
Versionamento & Arquivos finais estão nomeados com data, versão e responsável. & $\square$ \\
\bottomrule
\end{longtable}

% -------------------------------------------------
\section{Riscos e Controles}
% -------------------------------------------------

\begin{longtable}{p{0.30\textwidth}p{0.35\textwidth}p{0.25\textwidth}}
\toprule
\textbf{Risco} & \textbf{Impacto} & \textbf{Controle sugerido} \\
\midrule
Alteração indevida do JSON estrutural & Curso pode deixar de abrir ou perder blocos. & Editar apenas campos textuais e manter backup. \\
Prompt inadequado na geração & Curso raso, genérico ou desalinhado ao plano. & Registrar e versionar prompts aprovados. \\
Verificação visual insuficiente & Problemas passam para o SCORM final. & Usar checklist mínimo por lição. \\
Perda de rastreabilidade & Dificuldade para auditar versões. & Nomear arquivos com curso, data e versão. \\
Erro na reserialização & SCORM final pode ficar inválido. & Testar pacote antes da publicação. \\
\bottomrule
\end{longtable}

% -------------------------------------------------
\section{Registro de Versões}
% -------------------------------------------------

\begin{table}[h!]
\centering
\begin{tabularx}{\textwidth}{p{0.16\textwidth}p{0.18\textwidth}Yp{0.22\textwidth}}
\toprule
\textbf{Versão} & \textbf{Data} & \textbf{Alteração} & \textbf{Responsável} \\
\midrule
1.0 & \today & Criação inicial do relatório do fluxo. & João Victor  \\
 & & & \\
 & & & \\
\bottomrule
\end{tabularx}
\caption{Histórico de versões do documento.}
\end{table}

% -------------------------------------------------
\section{Conclusão}
% -------------------------------------------------

O fluxo proposto permite organizar a produção de cursos digitais de maneira controlada, combinando automação, revisão visual, manipulação técnica do pacote SCORM e preservação da estrutura interna do Articulate Rise 360. A principal cautela do processo está na separação entre conteúdo editável e estrutura técnica, especialmente na etapa de alteração dos arquivos JSON.

\end{document}
